\todo{Lecture 23}
\begin{theorem}
Soit $A\in \mathbf{C}^{n\times n}$. Ils existent $P\in \mathbf{C}^{n\times n}$ inversible et $J\in \mathbf{C}^{n\times n}$ en forme normale de Jordan telles que 
$$
A=PJP^{-1}.
$$
\end{theorem}
\begin{lemma}
Soit $J\in \mathbf{C}^{n\times n}$ en forme normale de Jordan. Alors 
$$
J=\underbrace{ \begin{pmatrix}
\lambda_{1}I_{n_{1}} &  &  \\
 & \ddots &  \\
 &  & \lambda _{k}I_{n_{k}}
\end{pmatrix} }_{ D }+\underbrace{ \begin{pmatrix}
J(0,n_{1}) &  &  \\
 & \ddots &  \\
 &  & J(0,n_{k})
\end{pmatrix} }_{ N }.
$$
\end{lemma}
\begin{proof}
On a
$$
N^{i}=\begin{pmatrix}
J(0,n_{1}) &  &  \\
 & \ddots &  \\
 &  & J(0,n_{k})
\end{pmatrix}^{i}=\begin{pmatrix}
J(0,n_{1})^{i} &   &  \\
 & \ddots &  \\
 &  & J(0,n_{k})^{i}
\end{pmatrix}.
$$
Soit $j=\max\{n_{1},\dots,n_{k}\}$, alors $N^{j}=0$. On a que $D$ et $N$ commutent 
$$
\lambda _{i}I_{n_{i}}J(0,n_{i})=J(0,n_{i})\lambda _{i}I_{n_{i}} \Rightarrow DN=ND.
$$
\end{proof}
\begin{reminder}
$$
A=P(D+N)P^{-1}\Rightarrow e^{A}=Pe^{ D }e^{ N }P^{-1}
$$
ou $e^{ N }$ est une somme fini car $N$ est nilpotente.
\end{reminder}

\question{Quel est le lien entre $J$ et $m_{A}(x)$ de $A$ ?}

\begin{reminder}
$m_{A}(x)\in \mathbf{C}[x]$ est le polynôme minimal de $A$ (un polynome unitaire et de degre minimale tel que $m_{A}(A)=0\in \mathbf{C}^{n\times n}$).
\end{reminder}

Soit $p(x)\in \mathbf{C}[x]$. On a 
\begin{align*}
p(A) & =p(PJP^{-1}) \\
 & =Pp(J)P^{-1}
\end{align*}
alors $p(A)=0$ si et seulement si $p(J)=0$. 

\begin{example}
On a avec $\lambda_{1}\neq\lambda_{2}$ :
$$
A \sim \underbrace{ \begin{pmatrix}
\lambda_{1} & 1 \\
 & \lambda _{1} & 1 \\
 &  & \lambda _{1} &  \\
 &  &  & \lambda _{2} \\
 &  &  &  & \lambda _{2}
\end{pmatrix} }_{ J }.
$$
Le polynôme caractéristique de $J$ est 
$$
f_{J}(x)=(-1)^{5}(x-\lambda _{1})^{3}(x-\lambda_{2})^{2},
$$
donc le polynôme minimal est
$$
m_{J}(x)=(x-\lambda_{1})^{e_{1}}(x-\lambda_{2})^{e_{2}}.
$$
On a 
$$
(J-\lambda_{1}I)^{e_{1}}=\begin{pmatrix}
0 & 1 \\
 & 0 & 1 \\
 &  & 0 \\
 &  &  & \lambda_{2}-\lambda_{1} \\
 &  &  &  & \lambda_{2}-\lambda_{1}
\end{pmatrix}^{e_{1}}.
$$
L'exponent $e_{1}$ est tel que
$$
\begin{pmatrix}
0 & 1 \\
 & 0 & 1 \\
 &  & 0
\end{pmatrix}^{e_{1}}=0\in \mathbf{C}^{3\times3}.
$$
Donc $e_{1}=3$ et $e_{2}=1$.
\end{example}
\begin{lemma}
Soit $J\in \mathbf{C}^{n\times n}$ en forme normale de Jordan et soient $\lambda_{1},\dots,\lambda _{j}\in \mathbf{C}$ des valeurs propres distincts de $A\in \mathbf{C}^{n\times n}$. Alors le polynôme minimal de $J$ est
$$
m_{J}(x)=(x-\lambda_{1})^{e_{1}}\cdots(x-\lambda _{k})^{e_{k}}
$$
ou $e_{i}=\max\{ n_{i}:J(\lambda _{i},n_{i})\text{ est un bloc de }J \}$. 
\end{lemma}
\begin{proof}
Soit 
$$
J=\begin{pmatrix}
J(\lambda_{1},n_{1}) \\
 & J(\lambda_{1},n_{2}) \\
 &  & \ddots  \\
 &  &  & J(\lambda_{1},n_{k}) \\
 &  &  &  &  J(\lambda_{2},n_{k+1}) \\
 &  &  &  &  &
\end{pmatrix}
$$
Alors $e_{1}$ est tel que
$$
\begin{pmatrix}
J(0,n_{1})^{e_{1}} \\
 &  \ddots \\
 &  &  J(0,n_{k})^{e_{1}}
\end{pmatrix}=0
$$
et pareil pours les autres blocs.
\end{proof}
\begin{example}
Soit $A\in \mathbf{C}^{5\times5}$ et $f_{A}(x)=(-1)(x-2)^{4}(x-5)$. Alors
$$
J=\begin{pmatrix}
2 & * \\
 & 2  & *\\
 &  & 2  & *\\
 &  &  & 2 \\
 &  &  &  & 5
\end{pmatrix}.
$$
Ici on a quelques possibilités de former les blocs selon si les $*$ sont 1 ou 0. 
\end{example}
\begin{reminder}
Si $A\sim J$, alors
$$
\mathrm{rank}(A^{i})=\mathrm{rank}(J^{i})\Leftrightarrow \dim(\ker(A^{i}))=\dim(\ker(J^{i})).
$$
\end{reminder}
\begin{theorem}
Soit $A\in \mathbf{C}^{n\times n}$ et $\lambda_{1},\dots,\lambda _{k}\in \mathbf{C}$ les valeurs propres distinctes de $A$. Soit $d_{i}=\dim(\ker(A-\lambda_{1}I)^{i})$ pour $i=1,\dots,n$. Alors $d_{i}-d_{i-1}$ est le nombre de blocs de $J$ dont $\lambda _{1}$ est sur la diagonale et de longueur au moins $i$.
\end{theorem}
\begin{proof}
On a
$$
d_{i}=d_{i-1}+\text{\# de blocs de longueur }\ge i.
$$
Alors le nombre de blocs de longueur au moins $i$ est egal a $d_{i}-d_{i-1}$. Le nombre de blocs de longueur exactement $i$ est donc 
$$
(d_{i}-d_{i-1})-(d_{i+1}-d_{i})=2d_{i}-d_{i-1}-d_{i+1}.
$$
\end{proof}

	